\documentclass[11pt,a4paper,UTF8]{ctexart}
\usepackage[margin=2.2cm]{geometry}
\usepackage{amsmath,graphicx,booktabs}
\usepackage[section]{placeins}
\usepackage[round,authoryear]{natbib}
\usepackage[colorlinks=true,allcolors=blue]{hyperref}
\setlength{\emergencystretch}{3em}
\setlength{\parskip}{0.25em}
\title{不同语言模型具有不同的教育人格吗？基于行为预测的研究}
\author{匿名工作稿中文阅读版}\date{}
\begin{document}\maketitle

工作稿。已包含引言、相关工作、方法、测量开发与局限；独立确认结果尚未产生，本文没有预写实证摘要和结论。实验数量表示冻结设计，实际完成量以最终审计为准。

\section{引言}

面对同一道题和同一段学生作答，两个语言模型可能都知道正确答案，却以不同方式教学：一个直接给出结果并解释，另一个邀请学生继续推理，还有一个先回应学生表达的挫败感。这些差异引出本研究的核心问题：不同模型是否具有可重复的教育行为倾向，能够帮助我们预测它在新情境下如何回应？我们将这种待检验的行为规律称为“教育人格”。

辨认出模型的文本并不足以回答这个问题。篇幅、排版、题目要求和答案是否可得，都可能产生鲜明的模型差异。同时，稳定也不意味着模型在任何情境下都采取相同行动。如果一个模型经常在学生出错时增加解释，在学生表达挫败时回应困难，这种随情境变化的规律同样可以重复。我们需要区分三个层次：观察到差异，可靠地测到差异，以及通过差异预测新的行为。

教育语境使这个问题具有可观察的行为对象。是否揭示答案、是否邀请学生推理，都能从实际回答中编码；两者可以同时出现，也没有哪一种永远更好。回应学生表达的困难同样可以观察，而无需推断模型是否体验同理心。本文研究的是导师模型的行为倾向，与 PATS 等根据学生人格调整教学的研究对象不同 \citep{rooein2026pats}，也不据此推断教学质量或学习增益。

已有研究通过人格情境和行为探针识别模型差异 \citep{lee2025trait,liu2026bma}，教育领域也已比较模型教学动作组合，并从对话中学习导师差异 \citep{kucheria2025cima,lee2026tutorpersonas}。本文在此基础上检验默认倾向和学生状态响应的额外预测信息，与情境、模型在不同学科的习惯及简单训练输出风格比较，并在新回答生成前锁定预测。

全文围绕三个相互衔接的问题展开：

\begin{enumerate}
\item \textbf{可重复差异}：相同教育输入下，不同模型是否表现出可区分且跨重复请求复现的默认倾向？
\item \textbf{新来源预测}：默认倾向或模型特有的学生状态响应，能否预测独立来源新题中的行为？
\item \textbf{提示可塑性}：中性角色措辞和明确教学指令分别如何改变这些倾向？
\end{enumerate}

我们将已有教育输出档案用于探索，并另行开展前瞻性确认。档案中的百万余条原始记录涉及不同角色、覆盖范围和重复历史，必须先审计，不能全部当作独立的导师人格证据。测量开发试验与正式材料永久分开；用训练材料拟合画像、保存确认预测后，才生成确认回答。结论可以支持默认倾向、进一步支持条件响应规律，也可能只支持局部边界，不能由研究动机预先决定。

\section{相关工作}

\textbf{人格与情境行为。} TRAIT 将人格量表扩展为情境选择，并考察可区分性、一致性和提示作用 \citep{lee2025trait}。BMA 在多种互动语体中使用情境选择，也评估开放生成和内部激活控制；其框架已区分构建情境与评估情境 \citep{liu2026bma}。另有研究通过内部特征干预联系人格表征与情境行为 \citep{zhang2026triad}，或观察相同人格指令在不同对话中的表达 \citep{han2026contexts}。这些研究共同支持检查默认倾向与条件表达。本研究预测部署输出中的教育动作，不据此识别其内部机制。

\textbf{导师行为与教育人格。} CIMA 研究在语言学习对话中比较三种模型的动作组合和回答复杂度，生成提示同时要求回答与动作标签 \citep{kucheria2025cima}。描述不同模型的教学风格已有直接前作。导师人格操控研究从人类教学对话学习差异，并在留出对话中评估人格对齐 \citep{lee2026tutorpersonas}。本文研究多个现有部署的默认行为，以及受控学生线索和教学指令如何改变这些行为；在生成后编码可见回答，并检查不同评审面板的影响。这改变了测量流程，但不保证编码就是真值。

\textbf{教学质量与行为预测。} MRBench 和 MathTutorBench 提供教学能力评价 \citep{maurya2025mrbench,macina2025mathtutorbench}，PATS 根据学生人格调整教学 \citep{rooein2026pats}。本文的主要目标是特定条件下导师行为的发生概率。可预测性不能直接排列其教育价值；我们要检验默认与条件画像在新来源上的增益、与简单解释的比较，以及同题中性措辞和明确指令的效应。来源分离和提前锁定预测使这些比较可以审计，实质贡献仍取决于得到的证据。

\section{研究设计与检验}

本文将教育人格操作化为特定部署反复出现、能够帮助预测新教育输入下行为的默认倾向及条件响应规律。研究单位包含请求与返回模型名、服务方、观察时间和生成参数；五个配置是 MiniMax-M3、MiniMax-M2.7、GLM-5.3、DeepSeek-V4-Pro 与 Doubao-Seed-2.0-Lite。这不假设人类人格结构，也不要求跨情境排名不变。完整材料、编码与估计规则、协议修改记录见文末附录。

\begin{figure}[htbp]\centering
\includegraphics[width=\linewidth]{figures/study_design.pdf}
\caption{前瞻性设计示意，数字表示计划回答量。每道主实验题交叉五个部署、四种学生作答／情绪条件和两次请求。训练与主确认来源家族分开；RQ1、RQ2 使用 32 个标准角色确认来源，RQ3 在其中 16 个来源上增加角色臂。重复请求和提示变体不增加独立来源数。四个嵌套预测器另与只用训练数据的风格预测器及后述敏感性分析比较。}
\end{figure}

\subsection{探索、测量开发与独立确认}

百万记录档案用于探索。保留输入变体区分的 899,025 条成功去重记录涉及不同角色，当前适配器审计识别其中 68,799 条来自九类导师回应任务；历史提示与部署溯源并不完整。因此档案发现与前瞻性确认分开，160 条回答的测量开发试验也永久排除于正式材料。

正式训练有 32 题、30 个保守来源家族，确认有另外 32 题、32 个来源家族；每个分区在数学、科学、语言、逻辑与数据推理四领域各有八题。相关训练题在内部验证中成组，测量试验模板及数值替换题均排除。研究 agent 编写题目及学生消息，两位模型评审在冻结前审核全部 64 套材料。确认目标是这四领域内的新来源，不是新学科或课堂总体。

每题交叉错误尝试／正确未完成作答与明确平静／挫败表述。作答对比同时改变正确性和进展，情绪对比只操控文本线索。各主条件提供相同教师参考解并说明学生未见，减少答案可得性差异，但不能消除能力或指令遵循差异。每模型—输入单元请求两次，温度 0.7、输出上限 8,192 token。训练与标准角色确认各 1,280 条回答，16 个来源上的额外提示臂共 2,560 条；表达机会探针和跨阶段哨兵分别增加 160、40 条，合计计划导师生成 5,320 条。重复和变体不增加独立来源数。

\subsection{可观察行为与重复性（RQ1）}

三个试验选定的主要事件是：答案揭示，即明确给出最终结论而不论其正确性；推理邀请，即邀请学生完成尚未被导师立即作答的推理；情绪／困难承接，即超出泛泛礼貌表扬的明确识别。这些动作可以共存，不是独立潜在人格因子。另保留五个次要事件，不依正式结果更换主要终点。

MiniMax-M3、DeepSeek-V4-Pro 和 GLM-5.3 分别使用学生上下文、参考目标与原始回答行编码。评审不接收生成模型名、角色臂、重复编号或分区标签，但风格仍可能泄露身份。正例需要有效非空证据行；主要标签为三份结构有效编码的多数。证据可追踪及评审一致都不等于语义真值。重试仅修复完成或结构失败，不按事件标签或分歧挑选。一份缺失训练编码在锁定预测前采用单项扩大预算修订，具体范围和原失败保留于附录。

RQ1 报告模型发生率、学生状态分层画像，以及两次请求的正例、负例和总体一致率，保留来源聚类。常量行为也可能完全一致，因此需同时考察发生率与预测价值。

\subsection{提前锁定的新来源预测（RQ2）}

四个嵌套逻辑预测器依次纳入学科×学生状态、模型身份、模型×学科、模型×学生状态。新特征块仅用训练数据正交化和标准化，使已有块的 ridge 惩罚几何保持一致。正则强度由训练来源分组验证从 0.1、1、10、100 中选择，截距不惩罚。拟合参数、调参结果与确认概率在任何确认导师回答生成前锁定；不使用确认答案或标签拟合。另设只用训练对数词数和标记行比例画像的描述性风格竞争者。

三个主要事件各比较默认画像相对情境基线、学生状态交互相对模型—学科画像的 Brier 改善。损失先在来源内平均，再对 32 个确认来源等权平均。六项单侧来源配对检验使用 Holm 校正，同时报告绝对误差、全部增益方向与来源 bootstrap 区间。区间以已拟合画像和来源抽样近似为条件，不代表全部教育情境或模型家族的不确定性。

\subsection{提示效应与替代解释（RQ3）}

冻结哈希选定的 16 个来源对每个模型完整交叉学生状态、标准角色、两种中性角色改写，以及明确 ASK、EXPLAIN 指令。来源内配对发生率差描述平均位移及模型／状态异质性。中性措辞等效参考为 ±0.10；保留经过多重性调整的名义配对 t 区间，但额外有界来源检查在 16 个来源下精度不足，不能建立总体等效。不显著不等于稳定。明确指令效应不表示教学改善；答案揭示与推理邀请的联合发生率允许两者共存，不推断先后顺序。

确认生成前还规定八种评审面板、五种删除单一生成配置的分析和 200 次训练来源重抽样。数学子集与内容错误筛选评估已锁定预测，不重新拟合；标记回答时排除完整五模型比较单元。两位内容评审另接收十二条 agent 编写的诊断控制。它们检查测量、配置、训练材料及错误依赖，不能建立人工真值或能力的因果控制。精确规则、制定时间及覆盖要求均保留于附录。

\section{测量开发试验结果}

本节仅报告已完成并永久排除于正式确认的测量试验：160 条导师回答来自四个模板，获得 480 份完整有效评审。它不提供教育总体的事件发生率或确认性置信区间。三个主要行为中，答案揭示的多数正例为 129，推理邀请为 61，情绪／困难承接为 111；两两正例一致率分别为 98.04\%—99.23\%、89.43\%—92.19\%、97.32\%—98.63\%。这支持将它们作为主要行为终点，但不等于验证了三维人格结构。

相比之下，索取缺失事实只有一个多数正例，正例一致率为 0\%—40\%；无依据能力评价的正例一致率为 33.33\%—51.16\%。这些结果说明必须分别检查正例和负例一致，而不能依赖总体高一致率。下图中的线段是三种评审配对的范围，不是置信区间。

\begin{figure}[htbp]\centering
\includegraphics[width=\linewidth]{figures/measurement_pilot.pdf}
\caption{测量开发试验：160 条导师回答、480 份有效编码。每个点表示一种评审配对，线段表示配对之间的范围；蓝色为正式生成前选定的三个主要终点。编码一致不能当作语义真值。}
\end{figure}

测量试验中的模型默认画像呈现探索性预测信息，但加入学生状态交互未显示一致的额外增益。因此正式实验使用独立来源、领域内多题以及提前保存的预测，不能用试验数据代填确认结论。

\section{尚待独立确认的结果与讨论}

正式结果应依次回答三个研究问题，并与测量、来源和部署边界对应。训练数据或测量试验中的差异不能填入确认结果。尤其需要检验默认画像的增益是否依赖个别配置或评审、条件画像是否提供额外信息，并报告中性措辞的实际位移及等效判断的精度限制。未显示条件增益，只能约束所研究的两类学生线索与预测方法，不能推出模型完全不具备情境响应能力。

最后的贡献判断仍取决于这些证据。数据规模、调用成功、统计程序通过测试以及论文能编译，都不能代替这一判断。

\section{研究局限}

\textbf{构念范围。} 教育人格在本文中表示教育情境中观察到的行为。由于没有匹配的非教育交互对照，不能证明某种规律只属于教育领域。三个主要事件只覆盖少量动作，可以共同出现，不是经独立验证的潜在人格因子。行为更可预测也不等于更有帮助：答案揭示包含错误结论，推理邀请的适用性取决于学习者需要。本研究没有测量学习结果。

\textbf{材料与迁移。} 正式材料是 agent 编写的英文问题、脚本化学生消息和提供给导师的参考解，不是从课堂交互中概率抽样得到。四个领域和两类明确学生线索允许受控比较，但确认范围不包括多语言教学、长程对话、隐含学生状态或高级学科内容。参考解减少答案可得性差异，不能消除理解与指令遵循差异。按来源划分阻止指定题目家族跨区，不能证明相对于模型预训练的内容新颖性。

\textbf{测量与部署。} 自动内容审核和三评审编码不需要新增人工标注，但可能共享模型偏差。证据行检查确认位置，排除同家族评审检查一种依赖，两者都不能保证语义正确。少见或未出现的事件必须结合表达机会解释。五个部署也不是模型家族的随机样本；服务名称和短期重复请求不能证明权重或行为的长期稳定。跨阶段哨兵混合了生成波动、编码波动和可能的部署变化，不能单独识别其原因。

\textbf{统计推断与研究历史。} 确认材料包含 32 个来源组，不能因有数千回答就当作数千独立教育情境。来源区间依赖对类似于本题库的问题变化的近似；微小增益仍可能无法确定，16 个来源的保守检查不能建立总体等效。条件预测受所选学生线索、特征和正则化模型类别限制，增益弱不能否定所有情境响应。档案和试验已参与研究设计，一部分敏感性与解释规则在训练期间才固定；我们披露时间并保留失败测量。本地哈希和执行时间支持审计先后关系，但不是独立第三方预注册或不可篡改的时间证明。

\section{附录：完整方法与协议记录}

教育人格在本文中指特定部署反复出现的默认行为及条件响应规律。默认画像是在指定输入分布和标准导师角色下的行为概率；条件画像进一步描述概率随学生状态的变化。两种画像都需要接受独立来源预测检验。我们不假设人类人格量表可以直接套用于模型，不要求跨情境模型排名完全不变，也不把任何可识别的语言特征都视作人格。

研究单位是部署配置，包括请求与返回模型名、服务提供方、观察时间及生成参数。前瞻性研究包含 MiniMax-M3、MiniMax-M2.7、GLM-5.3、DeepSeek-V4-Pro 和 Doubao-Seed-2.0-Lite。测量试验中请求 GLM-5.2 实际全部返回 GLM-5.3，因此按返回配置报告，不与历史 5.2 合并。服务端名称不能证明权重始终不变。





\subsection{档案探索与来源}

对全部 39 个任务的当前适配器角色审计表明，在保留输入变体区分的 899,025 条去重记录中，354,308 条属于评分或判断已有作答，306,307 条属于考试或选项选择，68,799 条来自九类导师回应任务。角色分类有代码依据，但不能证明每个历史请求都使用当前适配器。对九类导师任务的 122 个预测文件与对应摘要逐一核验，全部与冻结档案哈希一致；其中 86,072 个每文件最新成功项目均未保存本审计检查的完整输入、生成输入哈希、返回模型名或单次生成时间字段。后续若恢复这些信息，必须区分原始保存与后来重建。详细口径见 research/54。

\subsection{材料与交叉设计}

正式材料含训练题 32 道、确认题 32 道。每个分区在数学、科学、语言、逻辑与数据推理四个领域各有 8 道题。训练材料按保守规则归为 30 个来源家族，相关的平均数题与单位转换题在内部验证中成组；确认材料属于另外 32 个来源家族。四个测量试验模板及其数值替换题排除在外。本研究检验已有领域内的新来源迁移，不声称推广到新学科或模型训练时从未见过的概念。

研究 agent 编写题目、参考解、学生错误尝试和正确但未完成的作答。两位模型评审在设计冻结前审核全部 64 套材料，共 128 份完整审核。各份审核的五项布尔检查均为肯定，但仍有 25 份附有问题说明；原文与 agent 复核结论一并保留。这是内容质量检查，不是人工金标准。

每题交叉两种学生作答状态——错误尝试、正确但未完成——以及两种明确情绪表述——平静、挫败，共四种条件。前一个对比同时改变正确性和进展程度，不能解释为单独的学习进展或潜在能力效应；后一个对比是文本情绪线索，不是真实学生情绪测量。每道主实验题都向导师提供相同的正确参考解，并说明学生尚未看到它。这减少答案可得性差异，但不能完全控制模型理解、注意或指令遵循能力。

训练阶段使用标准导师角色。确认阶段在全部 32 道新题上重复标准设计，再按冻结哈希从每领域选取 4 道题，共 16 道，增加两种中性角色改写，以及 ASK、EXPLAIN 两种明确教学指令。前者改变角色措辞，后者改变要求的教学行为；同题完整交叉全部角色臂与学生状态。

每个模型—输入单元发起两次请求，温度 0.7、输出预算 8,192 token。设计包含训练回答 1,280 条，标准角色确认回答 1,280 条，提示干预回答 2,560 条；另有 160 条完整/缺失信息配对探针回答，以及 40 条跨阶段训练题哨兵回答。提示变体与哨兵不增加独立确认来源数。保留请求顺序、时间、参数、输入哈希和失败历史，在服务并发上限内交错执行模型请求。

\subsection{行为测量}

测量开发试验确定三个主要行为终点：

\begin{itemize}
\item \textbf{答案揭示}：明确给出最终结论，包括错误结论或正确指出信息不足。因此衡量是否揭示结论，不衡量正确性。
\item \textbf{推理邀请}：邀请学生进行尚未由导师完成的推理，包括祈使句和间接建议；导师立即自行作答的修辞性问题不计入。
\item \textbf{情绪／困难承接}：明确识别情绪或困难；普通礼貌用语和泛泛表扬不计入。
\end{itemize}

三个终点不等于三个独立人格因子。另保留分步解释、学习者选择、认识限定、索取缺失事实、无依据能力评价五项次要终点，无论结果如何都不替换主要终点。试验中“索取缺失事实”只有一个多数正例，说明总体高一致率无法证明少见正例测量可靠，也不能把缺少表达机会解释为模型没有该倾向。

MiniMax-M3、DeepSeek-V4-Pro 和 GLM-5.3 分别编码每条可见导师回答。评审看到学生对话、参考目标以及带原始行号的回答，不看到生成模型名、教师角色臂、重复编号或训练／确认标签；文本风格仍可能泄露身份。正例必须引用存在且非空的原始行，负例不附证据行。定位成功只证明证据可追踪，不保证语义正确。

主要分析使用三位有效评审的多数标签，保留逐评审结果及面板敏感性。初版逐字引用协议在自然回答上失败，测量开发阶段改为行号证据；原失败记录保留。最终试验的 160 条导师回答获得 480 份有效编码，其中两项结构／完成失败按相同输入重试。正式实验同样只修复结构或完成失败，不按事件标签或评审分歧筛选重试。

训练编码曾中断，复用有效记录并完成原定两轮修复后，有 3,839／3,840 份有效编码。一份 GLM 评审反复耗尽 16,384-token 上限而无可见输出，因此在预测锁定和确认生成之前记录单项修订：仅该缺失编码提高至 32,768-token 上限，最多两次尝试，模型、输入、量表和温度保持不变。该项首次尝试成功，实际使用 3,511 completion tokens；这不能识别先前失败的原因。两种预算分别审计后，合并严格不重叠的有效编码，最终 1,280 条训练回答均有三份有效评审。原始中断、失败尝试和修订记录全部保留；扩展预算的一份编码明确属于原统一预算协议的偏离。

\subsection{先预测，再生成确认回答}

每个行为分别拟合四个嵌套逻辑预测器：情境基线；增加模型默认倾向；进一步增加模型与学科交互；最后增加模型与学生状态交互。新特征块仅在训练数据中对已有特征正交化并标准化，来源等权，截距不惩罚。正则强度从 0.1、1、10、100 中通过训练来源分组验证选择。训练标签为常量时，四个预测器使用同一个平滑常量概率。

拟合参数、调参结果、输入哈希与确认概率在任何确认回答生成之前保存并锁定。主要预测器不使用确认标签、确认回答长度或同题其他模型得分。另设只用训练回答的平均对数词数与标记行比例构成的风格竞争画像，在内部验证中重新估计。它只提供简单风格信息的竞争比较，不能当作对风格的因果控制。

三个主要事件各检验两个比较：默认画像相对情境基线的 Brier 误差改善，以及学生状态画像相对模型—学科画像的改善，共六项主要检验。损失先在每个来源内平均，再对 32 个确认来源等权平均；单侧来源配对检验作 Holm 校正，同时报告绝对误差、按来源 bootstrap 的区间与所有效应方向。这些区间以已拟合的训练画像为条件，不能推广为所有模型和所有教育情境的不确定性。

\subsection{重复、提示变化与稳健性}

重复分析分别报告正例、负例和总体一致率。恒定行为本身可能带来完全一致，因此同时展示事件发生率。提示效应在 16 个完整交叉来源中配对计算，比较各模型在替代角色与标准角色下的行为概率。

中性措辞的等效容许范围事先固定为事件概率的 ±0.10。保留针对 30 个模型—事件—中性角色比较的名义 Bonferroni 配对 t 区间；没有显著差异不等于等效。ASK 与 EXPLAIN 衡量教学指令造成的行为位移，不衡量教学改善。答案揭示和推理邀请可以组合为四种状态，证据行号不用于推断行动先后。

若所有观测来源的差值恰好为常量，配对区间可能退化为零宽。保留冻结程序的数值输出，但标记这种退化，不用它建立总体等效结论。这项解释规则在代码检查期间、确认生成之前明确，不能冒称原始设计就已完整规定。

训练期间的进一步合成诊断显示，少见行为在非零来源方差下也可能出现 t 区间覆盖不足。因此在确认生成前又固定有界来源检查，要求它同时支持才作总体等效陈述。每个来源差值位于 [-1,1]；缩放 Hoeffding 界并考虑双侧及 30 个比较后，半宽为 \texttt{sqrt(2 * log(2 * 30 / 0.05) / n)} \citep{hoeffding1963}。n=16 时约为 0.941，无法支持 ±0.10 的容许范围。可以报告这批来源中观察到的位移与名义区间，但不据此宣称总体等效。这项保守检查仍依赖来源独立的抽样近似，不能使 agent 编写的题库成为代表性样本；原六项主要预测检验不变。

确认生成前另锁定八种评审面板与五种删除单一生成配置的预测。前者检查标签解释的依赖性，后者检查某个配置是否驱动整体比较；它们重复使用相同回答，不构成外部复制，也不是预测未知模型。训练来源另做 200 次整组重抽样，保留重复抽中来源的次数权重和原选正则化比例。这些敏感性在测量试验之后、训练评审执行中制定，不替换六项主要比较，也不冒称整个研究开始前的预注册。

\subsection{数学子集与内容错误敏感性}

协议还要求报告数学子集与明显内容错误的敏感性。具体实现规则在训练评审期间、确认生成之前固定。先在全部八个数学来源、320 条标准角色回答上评估原先锁定的预测，不重新拟合。

MiniMax-M3 与 DeepSeek-V4-Pro 另行检查全部 1,280 条标准角色确认回答，根据学生输入、参考解和原始回答行号，区分明确内容错误、未发现明确错误、无法确定。不给答案、提问或正确引述并纠正错误本身不计为错误。明确错误必须给出类别及有效原始行证据。

另有 12 条 agent 编写的控制回答，共增加 2,584 个计划评审请求。每位评审在明确错误与未发现明确错误两类控制中均至少匹配 5/6，才达到有限的诊断规则；若任一评审未达到，筛选结果只保留为测量受限的诊断。即使达标，也不能据此证明自然回答的所有错误都已被发现。原始导师生成与主要检验不增加。

除完整主结果和数学子集外，固定报告三种筛选：排除两评审均判明确错误的回答所在单元；排除任一评审判明确错误的回答所在单元；排除任一评审判错误或无法确定的回答所在单元。单元由来源、学生状态和重复编号确定，包含五个模型；触发时整体排除，使模型比较保持配对。报告保留覆盖、来源等权损失和描述性来源区间，不增加假设检验。这些筛选依赖生成后的回答并改变目标分布，不是对能力的因果控制；内容评审范围也不包含全部提示干预回答。

\bibliographystyle{plainnat}
\bibliography{references}
\end{document}
